\chapter{Второе задание.}
\section{Банаховы алгебры.}
\begin{task}[T1]
    Пусть $\A$ -- банахова алгебра, $x \in \A$ и $y \in \A$.
    \begin{enumerate}
        \item Доказать, что если $x$ и $xy$ обратимы в $\A$, то обратим и $y$.
        \item Доказать, что если $xy$ и $yx$ обратимы в $\A$, то обратимы $x$ и $y$.
        \item Показать на примере, что может быть $xy = e \neq yx$.
        \item Если $\dim \A < \infty$, то $yx = e$ эквивалентно $xy = e$.
        \item Доказать, что элемент $e - yx$ обратим, если обратим элемент $e - xy$.
        \item Пусть $\lambda \in \sigma(xy)$ и $\lambda \neq 0$.
        Доказать, что $\lambda \in \sigma(yx)$.
        Показать, что $\sigma(xy)$ может содержать ноль, тогда как $\sigma(yx)$ нуля не содержит.
        \item Доказать, что если $x$ -- обратим, то $\sigma(xy) = \sigma(yx)$.
        \item Доказать, что $xy$ и $yx$ имеют один и тот же спектральный радиус.
        \item Доказать, что если $P(x) = 0$ для многочлена из $\Cm[x]$ такого, что $P(0) \neq 0$, то $x$ -- обратим.
    \end{enumerate}
\end{task}
\begin{solution}
    \begin{enumerate}
        \item Рассмотрим $z = (xy)^{-1}x$. Очевидно, что $zy = e$, а значит $z$ -- левый обратный $y$. Но и $yz = y (xy)^{-1}x = x^{-1}xy(xy)^{-1}x = x^{-1}ex = e$, а значит $z$ -- и правый обратный. Поэтому $y^{-1} = z$ и $y$ -- обратим в $\A$.
        \item Заметим, что левый и правый обратный элемента $z \in \A$ если существуют, то совпадают.
        И действительно: пусть $u$ -- левый обратный $z$, а $v$ -- его правый обратный.
        Тогда $e = uz \Rightarrow v = uzv = u(zv) = ue = u$.
        Следовательно, если у элемента банаховой алгебры существуют левый и правый обратный, то они совпадают и элемент обратим.
        Но, по условию, $(xy)^{-1}xy = e$, то есть $((xy)^{-1}x)y = e$ -- то есть $y$ обратим слева.
        А с другой стороны $yx(yx)^{-1} = y(x(yx)^{-1}) = e$, а значит $y$ обратим и справа.
        Поэтому он обратим. Аналогично, $x$ -- обратим.
        \item Рассмотрим $\ell^2$, банахову алгебру $\A = \LL(\ell^2, \ell^2)$ и операторы $x$ -- левый сдвиг, $y$ -- правый сдвиг.
        Очевидно, что $xy = e$. Также, очевидно, что $yx \neq e$, поскольку $\ker yx \neq 0$.
        \item Рассмотрим операторы левого умножения $\forall a \in \A$, которые действуют как
        \[
            L_a(x) = ax
        \]  
        Очевидно, что $L_x$ имеет нулевое ядро. Пусть $L_x z = 0$.
        Тогда $z = yxz = y0 = 0$.
        Поэтому, в силу конечномерности, $\dim \im \A = \dim \ker L_x + \dim \im L_x = \dim \im L_x$ и оператор $L_x$ -- биективен.
        Поэтому существует $z \in \A$ такой, что $xz = L_x z = e$.
        А значит $x$ -- обратим справа.
        Но он обратим и слева, а значит левый и правый обратные совпадают и $xy = yx = e$.
        \item Пусть $e - xy$ -- обратим.
        Тогда рассмотрим элемент $z = y(e - xy)^{-1}x + e$.
        Рассмотрим $(e - yx)z\colon$
        \[
            (e - yx)\left(y(e - xy)^{-1}x + e\right) = y(e - xy)(e - xy)^{-1}x + (e - yx) = e,
        \]
        и $z(e - yx)\colon$
        \[
            \left(y(e - xy)^{-1}x + e\right)(e - yx) = y(e - xy)^{-1}(e - xy)x + (e - yx) = e.
        \]
        Что и даёт искомое утверждение.
        \item Пусть $\lambda \in \sigma(xy)$ и $\lambda \neq 0$. Тогда $\lambda e - xy$ -- не является обратимым. С другой стороны, если $\lambda \notin \sigma(yx)$, то $\lambda e - yx$ -- обратим. Поскольку $\lambda \neq 0$, то это эквивалетно обратимости $e - \frac{1}{\lambda}yx$. По предыдущему пункту это означает обратимость $e - x \frac{1}{\lambda }y$, что эквивалентно обратимости $\lambda e - xy$ -- противоречие. 
        %Примером $x, y$ и $\A$, в которых одна композиция содержит ноль, а вторая -- нет являются сдвиги в $\ell^2$.
        \item Заметим, что если $x$ -- обратим, то $xy = x(yx)x^{-1}$, а значит
        \[
            \lambda e - xy = x( \lambda e - yx)x^{-1}.
        \]
        Поскольку умножения на $x$ и $x^{-1}$ -- обратимые операторы, то $\lambda e - xy$ обратим тогда и только тогда, когда обратим $\lambda e - yx$, поэтому и спектры совпадают.
        \item По позапрошлому пункта для произвольных $x$ и $y$ верно $\sigma(xy) \setminus\{0\} = \sigma(yx) \setminus\{0\}$, а значит если в спектре есть точки помимо $0$, то они есть в обоих спектрах и, таким образом, $r(xy) = r(yx)$.
        \item Многочлен $P(x)$ записывается как
        \[
            P(x) = a_0e + a_1x + \ldots a_n x^n.
        \]
        Так как $P(x) = 0$, то
        \[
            -a_0 e = x(a_1 + \ldots a_n x^{n - 1}).
        \]
        Из того, что $P(0) \neq 0$ следует то, что $a_0 \neq 0$.
        Поэтому 
        \[
            e = x\left(\frac{a_1}{-a_0} + \ldots \frac{a_n}{-a_0}x^{n - 1}\right) = x q(x).
        \]
        Заметим, что $x$ и $q(x)$ -- коммутируют, а значит $xq(x) = q(x)x = e$ и $x$ -- обратим.
    \end{enumerate}
\end{solution}
\begin{task}[T2]
    Пусть $\A$ -- банахова алгебра, и $\G(\A)$ -- группа обратимых элементов $\A$.
    \begin{enumerate}
        \item Доказать, что множество $\G(\A)$ -- открыто в $\A$.
        \item Доказать, что для любой граничной точки $x$ множества $G(\A)$ и для любой последовательности $x_n \in \G(\A)$, сходящейся к $x$ выполнено $\|x_n^{-1}\| \rightarrow \infty$.
    \end{enumerate}
\end{task}
\begin{solution}
    Заметим, что если $a \in \A$ такой, что $\|a\| < 1$, то элемент $e - a$ -- обратим и 
    \[
        (e - a)^{-1} = \sum_{k = 0}^{\infty} a^k.
    \]

    Покажем открытость $\G(\A)$.
    Пусть $a \in \G(\A)$.
    Рассмотрим $B_{\delta}(a) \subset \A$, где $\delta = \frac{1}{2\|a^{-1}\|}$.
    Покажем, что $a - \delta a$, где $\delta a \in B_{\delta}(a)$ -- обратим.
    Его обратимость эквивалентна обратимости элемента $e - a^{-1}\delta a$.
    В силу предыдущего замечания, так как $\|a^{-1}\delta a\| \leq \|a^{-1}\|\|\delta a\| < \frac{1}{2}$, то это действительно так и, а значит $\G(\A)$ -- открыто.

    Пусть $a \in \partial \G(\A)$ и $\{a_n\} \subset \G(\A)$ такая, что $a_n \rightarrow a$.
    Предположим, что $\|a_n^{-1}\|$ не стремится к $\infty$.
    Значит, существует подпоследовательность такая, что $\|a_{n_k}^{-1}\| \leq C$ при некоторой $C > 0$. Так как $a_{n_k} \rightarrow a$, то при $k \geq K$ верно, что $\|a - a_{n_k}\| \leq \frac{1}{2C}$. Тогда, так как $a = a_{n_k} + (a - a_{n_k}) = a_{n_k}(e + a_{n_k}^{-1}(a - a_{n_k}))$, и $\|a_{n_k}^{-1}(a - a_{n_k})\| < \frac{1}{2}$, то $a$ -- обратим. 
    Противоречие.
\end{solution}
\begin{task}[T3]
    Пусть $\A$ -- банахова алгебра такая, что
    \[
        \exists M > 0 \ \forall x, y \in \A \hookrightarrow \|x\|\|y\| \leq M\|xy\|.
    \]
    Показать, что $\A$ изометрически изоморфна $\Cm$.
\end{task}
\begin{solution}
    Для всяких элементов банаховой алгебры $x$ и $y$ верно, что
    \[
        \|xy\| \leq \|x\|\|y\|.
    \]
    А значит 
    \[
        \|xy\| \leq \|x\|\|y\| \leq M\|xy\|.
    \]

    Пусть $z \in \partial \G(\A)$ и $z \neq 0$.
    Тогда существует последовательность $z_n \subset \G(\A)$ такая, что $z_n \rightarrow z$.
    В частности, для всякого $n \in \N\colon$
    \[
        \|z_n z_n^{-1}\| = 1 \leq \|z_n\|\|z_n^{-1}\| \leq M = M\|z_nz_n^{-1}\|.
    \]
    Заметим, что начиная с некоторого номера $N \in \N$ последовательность $\|z_n\|$ -- ограничена и отделена от нуля (так как сходится к $\|z\|$).
    Но, тогда и последовательность $\|z_n^{-1}\|$ является ограниченной -- противоречие с первой задачей. А значит граница группы обратимых элементов не содержит ненулевых элементов.
    В силу её открытости это возможно лишь только если всякий ненулевой элемент является обратимым (так как $\A \setminus \{0\}$ -- связно).

    Заметим, что спектр всякого элемента $x \in \A$ -- непуст.
    Тогда существует $\lambda_x$ такой, что $x - \lambda_x e$ -- необратим.
    Но есть единственный необратимый элемент -- $0$, а значит $x = \lambda_x e$.
    Рассмотрим отображение $\mathrm{eval}\colon \A \to \Cm, \ x \mapsto \lambda_x$.
    Очевидно, что это биективный гомоморфизм алгебр.
    Более того, так как
    \[
        \|\mathrm{eval(x)}\| = \|\lambda_x\| = \|\lambda_x e\| = \|x\|,
    \]
    а значит он изометричен.
    Следовательно $\A \cong \Cm$.
\end{solution}
\begin{task}[T4]
    Пусть $\A$ -- банахова алгебра, $x \in \A$.
    Доказать, что для любого открытого множества $V \subset \Cm$ вида $\sigma(x) \subset V$ существует $\delta > 0\colon$
    \[
        \forall y \in \A\colon \|y - x\| < \delta \hookrightarrow \sigma(y) \subset V.
    \]
\end{task}
\begin{solution}
    Пусть $R_x(\lambda) = (x - \lambda e)^{-1}$ и $\rho(x) = \Cm \setminus \sigma(x)$.
    Поскольку $\sigma(x) \subset V$, то $\Cm \setminus V \subset \rho(x)$.
    Расссмотрим $K = (\Cm \setminus V) \cap B_{\|x\| + 1}[0]$.
    Так как резольвента непрерывна на резольвентном множестве, и $K \subset \rho(x)$, то резольвента непрерывна на компакте $K$ и, следовательно, достигается $M = \sup_{K}\|R_x(\lambda)\|$.
    
    Пусть $\delta < \min\{1, \frac{1}{M}\}$.
    Покажем, что если $\|y - x\| < \delta$, то $\sigma(y) \subset V$.
    Достаточно показать что для всякой $\lambda \in \Cm \setminus V$ элемент $y - \lambda e$ -- обратим.
    Пусть, сначала, $\lambda \in K$.
    Распишем $y - \lambda e$ как
    \[
        y - \lambda e = (y - x) + (x - \lambda e) = (x - \lambda e)((x - \lambda e)^{-1}(y - x) + e).
    \]
    Достаточно показать, что $\|(x - \lambda e)^{-1}(y - x)\| < 1$, так как в таком случае второй множитель будет обратим в силу сходимости ряда Неймана.
    Но
    \[
        \|(x - \lambda e)^{-1}(y - x)\| = \|R_x(\lambda)(y - x)\| \leq M\delta < 1.
    \]
    Пусть теперь $\lambda \notin V$ и $|\lambda| > \|x\| + 1$.
    Так как
    \[
        \|y\| \leq \|x\| + \|x - y\| < \|x\| + 1.
    \]
    А значит $|\lambda| > \|y\|$ -- но тогда $\lambda\left(\frac{y}{\lambda} - e\right)$ обратим в силу сходимости ряда Неймана и доказательство завершено.
\end{solution}
\section{Спектр и резольвента.}
\begin{task}[7.1]
    Пусть $E$ -- банахово пространство, $A \in \LL(E)$.
    Показать, что $\sigma(A^n) = \{\lambda^n \mid \lambda \in \sigma(A)\}$.
\end{task}
\begin{solution}
    По определению
    \[
        \mu \in \sigma(A^n) \Longleftrightarrow \not\exists (A^n - \mu I)^{-1} \in \LL(E).
    \]
    Рассмотрим $\lambda_1, \ldots, \lambda_n$ -- все корни уравнения $z^n = \mu$.
    Тогда оператор $A^n - \mu I$ раскладывается как
    \[
        A^n - \mu I = \prod_{i = 1}^{n} (A - \lambda_i I).
    \]
    \begin{lemma}
        Пусть $\A$ -- банахова алгебра с единицей, $B_1, \ldots, B_n$ -- попарно коммутирующие элементы.
        Тогда следующие утверждения эквивалентны:
        \begin{enumerate}
            \item $\forall j$ элемент $B_j$ -- обратим.
            \item $\prod_{i = 1}^{n}B_i$ -- обратим.
        \end{enumerate}
    \end{lemma}
    \begin{proof}
        Очевидно, что если все элементы обратимы, то обратимо и их произведение.
        Пусть теперь обратимо произведение $P = \prod_{i = 1}^{n}B_i$.
        Тогда элемент $C_j\colon$
        \[
            C_j = \prod_{k \neq j}B_k \cdot P^{-1},
        \]
        является обратным элементом для $B_j$.
        И действительно, он является правым обратным, так как
        \[
            B_jC_j = B_j\prod_{k \neq j}B_k \cdot P^{-1} = PP^{-1} = I,
        \]
        так как все $B_j$ попарно коммутируют.
        
        Заметим, что $B_j$ коммутирует с $P^{-1}$.
        И действительно, так как $B_j$ коммутирует с $P$:
        \[
            PB_j = B_jP,
        \]
        то, домножая слева на $P^{-1}$ и справа на $P^{-1}$ получим:
        \[
            B_j P^{-1} = P^{-1}B_j.
        \]
        А значит левый обратный является и правым обратным, то есть в принципе честным обратным.
    \end{proof}
    
    Для алгебры операторов в нашем конкретном случае эта лемма применима, поскольку полиномы от $A$ коммутируют друг с другом.
    Таким образом, необратимость произведения эквивалентна необратимости хотя бы одного множителя из $A - \lambda_k I$ -- то есть принадлежности хотя бы одного из $\lambda_k$ спектру $A$.
\end{solution}
\begin{task}[7.7]
    В пространстве $C[0, 1]$ рассмотрим оператор:
    \[
        (Ax)(t) = \int_0^{t} x(s)ds.
    \]
    Найти его спектр и резольвенту.
\end{task}
\begin{solution}
    Сначала найдём $A^n$.
    Покажем по индукции, что
    \[
        A^n x(t) = \int_0^t \frac{(t - s)^{n - 1}}{(n - 1)!}x(s)ds.
    \]
    При $n = 1$ это, очевидно, выполнено.
    Найдём $A^{n + 1}\colon$
    \begin{multline*}
        A(A^n x)(t) = \int_0^t A^n x(s)ds = \int_0^t \int_0^s \frac{(s - r)^{n - 1}}{(n - 1)!}x(r)drds = \\ = \int_0^t x(r) \int_r^t \frac{(s - r)^{n - 1}}{(n - 1)!}dsdr = \int_0^t \frac{(s - r)^n}{n!}x(r)dr,
    \end{multline*}
    где перемена порядка интегрирования обосновывается при помощи теоремы Фубини (а у нас всё интегрируемо, ибо ограничено на пространстве конечной меры).

    Теперь оценим норму $A^n$.
    Так как
    \[
        \|A^n x\|_\infty \leq \frac{\|x\|_\infty}{(n - 1)!} \sup_{t \in [0, 1]} \int_0^t (t - s)^{n - 1}ds = \frac{1}{n!}\|x\|_\infty.
    \]
    Следовательно, спектральный радиус $A$ равен нулю (так как $r(A) \leq \lim \frac{1}{\sqrt[n]{n!}} \rightarrow 0$). Поэтому единственной точкой спектра $A$ является $0$.

    Найдём резольвенту при $\lambda \neq 0$.
    \[
        \int_0^t x(s)ds - \lambda x(t) = y(t).
    \]
    Пусть $z(t) = \int_0^t x(s)ds$.
    Тогда у нас поставлена задача Коши с начальным условием $z(0) = 0\colon$
    \[
        z'(t) - \frac{1}{\lambda}z(t) = -\frac{1}{\lambda}y(t).
    \]
    Решением однодного уравнения является $C e^{\frac{t}{\lambda}}$.
    Воспользуемся методом вариации постоянной:
    \[
        C' + \frac{1}{\lambda}Ce^{\frac{t}{\lambda}} - \frac{1}{\lambda}Ce^{\frac{t}{\lambda}} = -\frac{1}{\lambda}y(t) \Rightarrow C' = -\frac{1}{\lambda}y(t)e^{-\frac{t}{\lambda}}.
    \]
    Следовательно
    \[
        C(t) =  -\frac{1}{\lambda}\int_0^{t} y(s) e^{-\frac{s}{\lambda}}ds,
    \]
    поэтому 
    \[
        z(t) = -\frac{1}{\lambda}\int_{0}^t e^{\frac{t - s}{\lambda}}y(s)ds.
    \]
    Отсюда выражается $x$ и сама резольвента:
    \[
        (A - \lambda I)^{-1} = -\frac{1}{\lambda}y - \frac{1}{\lambda^2}\int_0^t e^{\frac{t - s}{\lambda}}y(s)ds.
    \]
\end{solution}
\begin{task}[7.11]
    В пространстве $C[0, 2\pi]$ рассмотрим оператор
    \[
        (Ax)(t) = e^{it}x(t).
    \]
    Показать, что спектр $A$ -- единичная окружность, и ни одна точка спектра не является собственным числом.
\end{task}
\begin{solution}
    Очевидно, что если $\lambda\colon |\lambda| \neq 1$, то $A - \lambda I$ непрерывно обратим и обратный оператор задаётся формулой
    \[
        (A - \lambda I)^{-1}y = \frac{y}{e^{it} - \lambda}.
    \]
    Так как $S^1$ -- замкнута и точка замкнута, то они отделимы друг от друга, и $\frac{y}{e^{it} - \lambda}$ будет являться корректно определённой непрерывной функцией.

    Если же $|\lambda| = 1$, то $\lambda = e^{i\alpha}$.
    Тогда $(A - \lambda I)x(\alpha) = 0$.
    Так как существуют непрерывные функции, не равные нулю на $\alpha$, то оператор не сюръективен и, следовательно, необратим -- а значит $\lambda$ является точкой спектра.

    Предположим, что $\lambda \in \sigma(A)$ -- собственное число.
    То есть существует некоторая непрерывная функция $x$ такая, что $Ax = e^{i\alpha} x$.
    Но $Ax = e^{it}x$, и уравнение принимает вид $(e^{i\alpha} - e^{it})x(t) = 0$. 
    Таким образом, $x$ должна быть равна нулю везде кроме, быть может, одной точки. Тогда в силу продолжения по непрерывности $x$ на всю область определения мы получаем что $x = 0$.
    А значит $\lambda$ не является собственным числом.
\end{solution}
\begin{task}[7.13]
    Какие множества на комплексной плоскости могут являться спектром некоторого ограниченного оператора в $l^2$?
\end{task}
\begin{solution}
    Непустые компакты.
    Пусть $K \subset \Cm$.
    В нём есть счётное всюду плотное множество точек $\{\lambda_k\}$.
    Рассмотрим оператор $A\colon \ell^2 \to \ell^2$, действуюший как 
    \[
        Ax(k) = \lambda_k x(k).
    \]
    Он, очевидно, является ограниченным, поскольку $\{\lambda_k\} \subset K$ -- а значит $|\lambda_k| \leq C$ и, таким образом, $\|Ax\| \leq C\|x\|$.
    Для определённого таким образом оператора $\{\lambda_k\} \subset \sigma(A)$.

    Пусть $\lambda \notin K$
    Тогда $\rho(\lambda, \{\lambda_k\}) = d > 0$.
    Рассмотрим $A - \lambda I\colon$
    \[
        (A - \lambda I)x = y.
    \]  
    Или, покоординатно:
    \[
        (\lambda_k - \lambda)x(k) = y(k) \Rightarrow x(k) = \frac{y(k)}{\lambda_k - \lambda}.
    \]
    То есть у нас определён оператор $(A - \lambda I)^{-1}$.
    Покажем его непрерывность:
    \[
        \|(A - \lambda I)^{-1}y\|^2 = \sum_{k = 1}^{\infty} \frac{|y(k|^2}{|\lambda_k - \lambda|^2} \leq \sum_{k = 1}^{\infty} \frac{|y(k|^2}{d^2} = \frac{1}{d^2}\|y\|^2.
    \]
    Таким образом $\lambda \notin \sigma(A)$, и $\sigma(A) \subset K$.
    Но, спектр всегда замкнут, и в него вложена последовательность $\{\lambda_k\}$.
    Следовательно, в него вложено и замыкание этой последовательности.
    В силу её плотности в $K$ мы и получаем необходимое равенство $\sigma(A) = K$.
\end{solution}
\begin{task}[7.15]
    Найти спектр оператора $A \in \LL(L^2(\R))\colon$
    \[
        (Af)(x) = \int_{\R} \frac{f(y)dy}{1 + (x - y)^2}.
    \]
\end{task}
\begin{solution}
    Рассмотрим $\kappa(x) = \frac{1}{1 + x^2}$. Очевидно, что $\kappa \in L^2(\R) \cap L^1(\R)$.
    Оператор $A$ выражается как 
    \[
        Af(x) = \kappa * f(x).
    \]
    Таким образом, $A$ унитарно эквивалентен оператору умножения на $\mathcal{F}[\kappa](x) = \pi e^{-|x|}$, а значит \[\sigma(A) = \essim \mathcal{F}[\kappa] = [0, \pi].\]
\end{solution}
\begin{task}[11.9]
    Пусть $A \in  \LL(\ell^2)$ -- оператор правого сдвига, то есть $Ax(n) = x(n - 1)$. Найти $\sigma(A)$ и $\sigma(A^*)$.
\end{task}
\begin{solution}
    Заметим, что $A$ -- изометричен.
    Тогда $\|A\| = 1$, а, следовательно, и $\|A^*\| = 1$.
    Поэтому $\sigma(A) \subset B_1(0)$ и $\sigma(A^*) \subset B_1(0)$.

    Известно, что $A^*$ -- левый сдвиг.
    Пусть $\lambda\colon |\lambda| < 1$.
    Рассмотрим $x_\lambda\colon x_\lambda(n) = \lambda^n$.
    Это корректно определённый элемент $\ell^2$, и $A^*x_\lambda = \lambda x_\lambda$, а значит всякий элемент открытого единичного круга лежит в спектре $A^*$.
    Следовательно, так как спектр -- замкнут, то и замыкание открытого единичного круга лежит в спектре, и, таким образом, $\sigma(A^*) = B_1[0]$.

    Так как для гильбертового сопряженённого оператора $\sigma(A) = \overline{\sigma(A^*)}$, то $\sigma(A) = B_1[0]$.
\end{solution}
\begin{task}
    Найти спектр и резольвенту оператора Вольтерра $A\colon L^2([0, 1]) \to L^2([0, 1])$ вида
    \[
        Af(x) = \int_0^x f(t)dt.
    \]
\end{task}
\begin{solution}
    Сначала найдём $A^n$.
    Покажем по индукции, что
    \[
        A^n x(t) = \int_0^t \frac{(t - s)^{n - 1}}{(n - 1)!}x(s)ds.
    \]
    При $n = 1$ это, очевидно, выполнено.
    Найдём $A^{n + 1}\colon$
    \begin{multline*}
        A(A^n x)(t) = \int_0^t A^n x(s)ds = \int_0^t \int_0^s \frac{(s - r)^{n - 1}}{(n - 1)!}x(r)drds = \\ = \int_0^t x(r) \int_r^t \frac{(s - r)^{n - 1}}{(n - 1)!}dsdr = \int_0^t \frac{(s - r)^n}{n!}x(r)dr,
    \end{multline*}
    где перемена порядка интегрирования обосновывается при помощи теоремы Фубини (а у нас всё интегрируемо в силу неравенств Юнга и Гёльдера).

    Теперь оценим норму $A^n$.
    Так как
    \[
        \|A^n x\|_2 \leq \|\frac{(s - r)^n}{n!}\|_1\|x\|_2 \leq \frac{\|x\|_2}{n!}.
    \]
    Следовательно, спектральный радиус $A$ равен нулю (так как $r(A) \leq \lim \frac{1}{\sqrt[n]{n!}} \rightarrow 0$). Поэтому единственной точкой спектра $A$ является $0$.

    Найдём резольвенту при $\lambda \neq 0$.
    \[
        \int_0^t x(s)ds - \lambda x(t) = y(t).
    \]
    Пусть $z(t) = \int_0^t x(s)ds$.
    Тогда у нас поставлена задача Коши с начальным условием $z(0) = 0\colon$
    \[
        z'(t) - \frac{1}{\lambda}z(t) = -\frac{1}{\lambda}y(t).
    \]
    Решением однодного уравнения является $C e^{\frac{t}{\lambda}}$.
    Воспользуемся методом вариации постоянной:
    \[
        C' + \frac{1}{\lambda}Ce^{\frac{t}{\lambda}} - \frac{1}{\lambda}Ce^{\frac{t}{\lambda}} = -\frac{1}{\lambda}y(t) \Rightarrow C' = -\frac{1}{\lambda}y(t)e^{-\frac{t}{\lambda}}.
    \]
    Следовательно
    \[
        C(t) =  -\frac{1}{\lambda}\int_0^{t} y(s) e^{-\frac{s}{\lambda}}ds,
    \]
    поэтому 
    \[
        z(t) = -\frac{1}{\lambda}\int_{0}^t e^{\frac{t - s}{\lambda}}y(s)ds.
    \]
    Отсюда выражается $x$ и сама резольвента:
    \[
        (A - \lambda I)^{-1} = -\frac{1}{\lambda}y - \frac{1}{\lambda^2}\int_0^t e^{\frac{t - s}{\lambda}}y(s)ds.
    \]
\end{solution}
\begin{task}
    Пусть число $a \neq 0$. Найти спектр и резольвенту оператора трансляции $T_a\colon L^2(\R) \to L^2(\R)$ вида $T_af(x) = f(x + a)$.
\end{task}
\begin{solution}
    Заметим, что оператор трансляции унитарно эквивалентен (при помощи преобразовани Фурье) оператору $m_a\colon L^2(\R) \to L^2(\R)$, действующему как оператор умножения на $e^{ita}$.
    Поэтому, спектром $T_a$ будет являться существенный образ $e^{ita}$ -- то есть единичная 
    окружность.
\end{solution}
\section{Эрмитово-самосопряжённые ограниченные операторы.}
\begin{task}[11.5]
    Пусть $A$ -- самосопряжённый оператор в гильбертовом пространстве $H$.
    Доказать, что 
    \begin{enumerate}
        \item $\|A\| = \sup\limits_{\|x\| \leq 1} |\langle Ax, x \rangle|$
        \item $\|A\| = \sup\limits_{\|x\| = 1, \|y\| = 1}|\langle Ax, y \rangle|$
    \end{enumerate}
\end{task}
\begin{solution}
    Сначала покажем второе равенство.
    Заметим, что $\forall x, y\colon \|x\| = \|y\| = 1$ верно
    \[
        |\langle Ax, y \rangle| \leq \|Ax\|\|y\| \leq \|A\|.
    \]
    С другой стороны, если взять $y = \frac{Ax}{\|Ax\|}$, то
    \[
        \left|\langle Ax, \frac{Ax}{\|Ax\|}\right| = \|Ax\|,
    \]
    а значит
    \[
        \sup\limits_{\|x\| = 1, \|y\| = 1} |\langle Ax, y \rangle| \geq  \sup\limits_{\|x\| = 1} \|Ax\| = \|A\|.
    \]

    Теперь покажем первое равенство.
    С одной стороны при $\|x\| \leq 1$
    \[
        |\langle Ax, x\rangle| \leq \|A\|\|x\|^2 \leq \|A\|,
    \]
    Положим $q(x) = \langle Ax, x \rangle$.
    Заметим, что в силу самосопряжённости $A$ форма $q$ -- вещественна.
    Пусть $x, y \in H$. Пусть $z \in H$ такое, что $\Re \langle Ax, z \rangle = |\langle Ax, y \rangle |$ (для этого достаточно домножить $y$ на фазу).
    При всяком $t > 0$
    \[
        q(x + tz) - q(x - tz) = 4t \Re \langle Ax, z \rangle
    \]
    С другой стороны, $q(x + tz) \leq \sup_{\|x\| \leq 1}|\langle Ax, x \rangle| \|x + tz\|^2$
    и $-q(x - tz) \leq \sup_{\|x\| \leq 1}|\langle Ax, x \rangle| \|x - tz\|^2 $.
    А значит (с учётом тождества параллелограмма)
    \[
        4t|\langle Ax, y \rangle| \leq 2\sup_{\|x\| \leq 1}|\langle Ax, x \rangle| \cdot (\|x\|^2 + t^2\|y\|^2).
    \]
    То есть
    \[
        |\langle Ax, y \rangle| \leq \frac{1}{2}\sup_{\|x\| \leq 1}|\langle Ax, x \rangle| \cdot \left(\frac{\|x\|^2}{t} + t \|y\|^2\right).
    \]
    Заметим, что минимум по $t$ правой стороны неравенства достигается на $t = \frac{\|x\|}{\|y\|}$.
    А значит для произвольных $x, y$ верно, что
    \[
        |\langle Ax, y \rangle| \leq \sup_{\|x\| \leq 1}|\langle Ax, x \rangle| \|x\|\|y\|.
    \]
    Переходя к $\|x\| = \|y\| = 1$ и супремуму по левой части с учётом второго утверждения задачи получаем оценку снизу, на чём доказательство первого утверждения и завершается.
\end{solution}
\begin{task}[11.7]
    Пусть $A$ -- самосопряжённый неотрицательный оператор в гильбертовом пространстве $H$.
    Показать, что оператор $(I + A)^{-1}$ -- существует.
\end{task}
\begin{solution}
    Пусть $x \in H$.
    Тогда
    \[
        \|(A + I)x\|^2 = \|x\|^2 + 2\langle x, Ax \rangle + \|Ax\|^2 \geq \|x\|^2 \geq 0.
    \]
    
    В силу положительности оператора это равно нулю если и только если $x = 0$ и $Ax = 0$.
    Следовательно, $\ker I + A = 0$.

    Покажем замкнутость $\im A + I$.
    Пусть $y_n = (A + I)x_n \rightarrow y$.
    Тогда
    \[
        \|x_n - x_m\| \leq \|(A + I)(x_n - x_m)\| \leq \|y_n - y_m\| \rightarrow 0,
    \]
    и $\{x_n\}$ -- фундаментальна. Следовательно, у неё существует предел и образ замнкут.

    А значит
    \[
        \im A + I = (\ker (A + I)^*)^{\perp} = \ker (A + I)^\perp = 0^\perp = H,
    \]
    а значит существует $(A + I)^{-1}$.
\end{solution}
\begin{task}[11.11]
    Пусть $H$ -- гильбертово пространство и $A\colon H \to H$ -- симметричный линейный оператор.
    Показать, что $A \in \LL(H)$.
\end{task}
\begin{solution}
    Достаточно показать замкнутость графика.
    Пусть $(x_n, Ax_n) \subset \graph A$ и $(x_n, Ax_n) \rightarrow (x, y)$.
    Заметим, что тогда
    \[
        \langle Ax_k, z \rangle \rightarrow \langle y, z\rangle.
    \]
    С другой стороны
    \[
        \langle Ax_k, z \rangle = \langle x_k, Az \rangle \rightarrow \langle x, Az \rangle = \langle Ax, z\rangle.
    \]
    Так как это верно для произвольного $z$, то $Ax = y$ и график замкнут, а значит и $A \in \LL(H)$.
\end{solution}
\begin{task}[11.13]
    Пусть $H$ -- гильбертово пространство, $A \in \LL(H)$ -- самосопряжённый оператор.
    Доказать, что $\|A^n\| = \|A\|^n$ при всяком $n \in \N$.
\end{task}
\begin{solution}
    Известно, что для произвольного ограниченного оператора $A$
    \[
        \|A^* A\| = \|A\|^2.
    \]
    Так как $A$ -- самосопряжён, то
    \[
        \|A^2\| = \|A\|^2.
    \]
    Очевидно, что это верно и для всех степеней двойки.
    Пусть $n \in \N$.
    С одной стороны, из субмультипликативности операторной нормы
    \[
        \|A^n\| \leq \|A\|^n.
    \]
    С другой стороны, если $n < 2^{k}$, то 
    \[
        \|A\|^{2^k} = \|A^{2^k}\| = \|A^n A^{2^k - n}\| \leq \|A^n\| \|A\|^{2^k - n},
    \]
    из чего и следует $\|A\|^n \leq \|A^n\|$ -- из чего и следует нужное утверждение.
\end{solution}
\begin{task}[12.18]
    Найти норму оператора Вольтерра 
    \[
        Af(x) = \int_0^x f(t)dt,
    \]
    в $\LL(L^2([0, 1]))$.
\end{task}
\begin{solution}
    Для начала найдём гильбертов сопряжённый оператор к $A$:
    \[
        \langle x, Ay\rangle = \int \overline{x}Aydt = \int_0^1 \overline{x}(t)\int_0^t y(s)dsdt = \int_0^1 \int_t^1 \overline{x}(t)dt y(s)ds = \langle A^* x, y \rangle,
    \]
    где перемена пределов интегрирования допустима в силу совместной интегрируемости $x$ и $y$ на квадрате ($\|xy\|_1 \leq \|1\|_2\|xy\|_2 \leq \|x\|_2\|y\|_2 < +\infty$).

    Таким образом 
    \[
        A^* x(t) = \int_t^1 x(s)ds.
    \]
    Следовательно
    \[
        A^* A x(t) = \int_t^1 Ax(s)ds = \int_t^1 \int_0^s x(r)drds = \int_0^1 (1 - \max\{t, r\})f(r)dr.
    \]
    Аналогично
    \[
        AA^*x(t) = \int_0^1 \min\{t, r\}f(r)dr.
    \]
    Заметим, что это интегральные операторы с $L^2$-ядрами (на $[0, 1]^2$) и, таким образом, они компактны. Поэтому их спектр состоит из собственных значений и, быть может, нуля.

    Пусть $u = A^* A f$.
    Тогда
    \[
        u'(x) = -\int_0^x f(s)ds, \quad u''(t) = -f(t).
    \]
    Если $A^* A f = \lambda f$, то $u = \lambda f = -\lambda u''$, и мы получаем 
    \[
        f'' + \frac{1}{\lambda} f = 0.
    \]
    В силу вида $u$ мы получаем, что $u(1) = 0, \ u'(0) = 0$, а значит и $f(1) = 0,\ f'(0) = 0$.
    Тогда, так как общее решение имеет вид
    \[
        f(t) = C \sin \left(\frac{t}{\sqrt{\lambda}}\right) + D \cos \left(\frac{t}{\sqrt{\lambda}}\right).
    \]
    Из граничных условий $D \cos \left(\frac{1}{\sqrt{\lambda}}\right) = 0$ и $C = 0$.
    Следовательно
    \[
        \lambda_n = \frac{1}{\left(\frac{\pi}{2} + \pi n\right)^2}.
    \]
    Так как $A^* A$ -- интегральный оператор с ядром из $L^2$, то он компактен и, при этом, он является самосопряжённым -- поэтому $\|A^* A\| = \frac{4}{\pi^2}$ -- наибольшему собственному значению. 
    А значит, поскольку $\|A\| = \sqrt{\|A^*A\|}$, то $\|A\| = \frac{2}{\pi}$.

    Заметим, что $\sigma(A^* A) = \{\lambda_n\} \cup \{0\}$, так как $\lambda_n \rightarrow 0$.
    Аналогично $\sigma(AA^*) = \sigma(A^* A) = \{\lambda_n\} \cup \{0\}$ -- так как оператор $A^* A$ является самосопряжённым.
    При этом, $0$, очевидно, не является собственным числом ни для $A^* A$, ни для $AA^*$, а значит он лежит в непрерывном спектре (остаточный спектр самосопряжённого оператора пуст).
\end{solution}
\begin{task}[T7]
    Пусть $H$ -- гильбертово пространство и $A \in \LL(H)$.
    Пусть $A^* A$ -- компактен.
    Показать, что $A$ -- компактен.
\end{task}
\begin{solution}
    Пусть $\{x_n\} \subset H$ -- ограниченная последовательность.
    Тогда $\{A^*A x_n\}$ -- вполне ограничено, и существует фундаментальная подпоследовательность $\{A^* A x_{n_k}\}$.
    Но, с другой стороны 
    \[
        \|Ax_{n_s} - Ax_{n_r}\|^2 = \langle A(x_{n_s} - x_{n_r}), A(x_{n_s} - x_{n_r}) \rangle = \langle A^* A (x_{n_s} - x_{n_r}), x_{n_s} - x_{n_r} \rangle \leq \|A^*A(x_{n_s} - x_{n_r})\| \cdot 2C,
    \]
    где $\|x_n\| \leq C$ -- так как $\{x_n\}$ ограничена.
    А значит и $\{Ax_{n_s}\}$ -- фундаментальна, и, следовательно, $AB_1(0)$ -- вполне ограничен.
\end{solution}
\begin{task}[T8]
    Доказать, что оператор инверсии $J\colon L^2(\R) \to L^2(\R)$ вида
    \[
        Jf(x) = f(-x),
    \]
    является самосопряжённым и найти его спектр и резольвенту.
\end{task}
\begin{solution}
    Найдём $J^*$.
    По определению гильбертового сопряжённого
    \[
        \langle J^* f, g\rangle = \langle f, Jg \rangle.
    \]
    То есть
    \[
        \int_{\R} \overline{J^* f}(x)g(x)dx = \int_{\R}\overline{f}(x)g(-x)dx.
    \]
    А значит
    \[
        \int_{\R} \overline{f}(-x)g(x)dx = \int_{\R} \overline{J^* f}(x)g(x)dx,
    \]
    то есть $J^* f(x) = f(-x) = J f(x)$ -- и, таким образом, $J$ -- самосопряжённый оператор.

    Исследуем его спектр.
    Заметим, что $J^2 = I$.
    Следовательно, так как $\sigma(J^2) = \sigma(J)^2$, то $\sigma(J)^2 = \{1\}$, а значит спектр $J$ состоит из не более чем двух точек -- $1$ и $-1$.
    Обе точки являются собственными числами, с собственным векторами в виде чётных и нечётных функций соответственно.

    Если же $\lambda \neq \pm 1$, то из $J^2 = I$ следует
    \[
        (J - \lambda I)(J + \lambda I) = (1 - \lambda^2)I \Rightarrow (J -\lambda I)^{-1} = \frac{J + \lambda I}{1 - \lambda^2}.
    \]
\end{solution}
\begin{task}[T9]
    Пусть $A\colon L^2([0, 1]) \to L^2([0, 1])$ -- оператор Вольтерра.
    Найти спектр и резольвенту $AA^*$ и $A^* A$.
\end{task}
\begin{solution}
    Спектр найден ранее.


\end{solution}
\begin{task}[T10]
    Пусть $h \in L^\infty([0, 1])$ и $A\colon L^2([0, 1]) \to L^2([0, 1])$ имеет вид
    \[
        (Af)(x) = h(x)f(x) \quad a.e.\  x \in [0, 1], \quad \forall f \in L^2([0, 1]).
    \]
    Найти норму и спектр оператора $A$, указав все компоненты спектра.
\end{task}
\begin{solution}
    Так как
    \[
        \|Af\|_2^2 = \int \overline{Af(x)}Af(x)dx = \int |h(x)|^2|f(x)|^2 dx \leq \|h\|_\infty^2 \|f\|_2^2,
    \]
    а значит $\|A\| \leq \|h\|_\infty$.
    Покажем и обратное неравенство.
    Пусть $E_\epsilon = \{x \in [0, 1] \mid |h(x)| \geq \|h\|_\infty - \epsilon\}$.
    Тогда $A$ достигает $\|h\|_\infty$ на последовательности \[f_n = \frac{\Ind_{E_{1/n}}}{\sqrt{\LL^1(E_{1/n})}}.\]

    Найдём спектр.
    Для $\lambda \in \Cm$ оператор $A - \lambda I$ имеет вид
    \[
        (A - \lambda I)f(x) = (h(x) - \lambda)f(x).
    \]
    Определим существенный образ $\essim h$ как 
    \[
        \essim h = \left\{\lambda \in \Cm \mid \forall \epsilon > 0 \hookrightarrow \LL^1(\{x \mid |h(x) - \lambda| < \epsilon\}) > 0\right\}.
    \]
    Покажем, что $\sigma(A) = \essim h$.
    Пусть $\lambda \notin \essim h$.
    Тогда существует $\epsilon > 0$ такой, что $\LL^1(\{x \mid |h(x) - \lambda| < \epsilon\}) = 0$.
    А значит a.e. $|h - \lambda| \geq \epsilon$ и, следовательно, $\frac{1}{h - \lambda} \in L^\infty([0, 1])$. Тогда $A - \lambda I$ непрерывно обратим и обратный задаётся как 
    \[
        (A - \lambda I)^{-1}f = \frac{f}{h - \lambda}.
    \]
    Таким образом $\sigma(A) \subset \essim h$.

    Пусть теперь $\lambda \in \essim h$.
    При всяком $n \in \N$ множество $E_n = \{x \mid |h(x) - \lambda | < \frac{1}{n}\}$ имеет положительную меру.
    Рассмотрим $f_n = \frac{\Ind_{E_n}}{\sqrt{\LL^1(E_n)}}$1.
    Тогда $\|f_n\|_2 = 1$ и
    \[
        \|(A - \lambda I)f_n\|_2^2 = \frac{1}{\LL^1(E_n)}\int_{E_n}|h(x) - \lambda|^2 dx \leq \frac{1}{n^2} \rightarrow 0.
    \]
    Если $(A - \lambda I)$ непрерывно обратим, то 
    \[
        1 = \|f_n\|_2 = \|(A - \lambda I)^{-1}(A - \lambda I)f_n\| \leq \|(A - \lambda I)^{-1}\|\|(A - \lambda I)f_n\| \rightarrow 0,
    \]
    что и приводит нас к $\lambda \in \sigma(A)$.

    Очевидно, что в точечном спектре лежат атомы p-f на $\Cm$, то есть
    \[
        \sigma_p = \{\lambda \mid \LL^1(\{x \mid h(x) = \lambda\} > 0)\}.
    \]
    И действительно, если $f$ -- собственный вектор, то
    \[
        (A - \lambda I)f = 0
    \]
    почти всюду.
    Если мера точек из $h^{-1}(\lambda)$ равна нулю, то $f$ почти всюду должна быть равна нулю -- чего не может быть у собственного вектора.

    Оставшиеся точки лежат в непрерывном спектре, поскольку сопряжённый оператор $(A - \lambda I)^*$ -- так же оператор умножения, причём множество $\{\overline{h}(x) = \overline{\lambda}\}$ имеет нулевую меру, а значит ядро -- тривиально и тогда ортогональное дополнение $[\im A - \lambda I]$ тривиально -- то есть $[\im A - \lambda I] = L^2([0, 1])$.
\end{solution}